%============================================================
% FORMAT FORMAT AKHIR — Universitas Amikom Purwokerto
% Semua aturan format naskah dikumpulkan di sini:
%   - paragraf & spasi
%   - format BAB / sub-bab (A., B., ...)
%   - daftar isi / tabel / gambar / lampiran
%   - nomor halaman
%   - caption tabel & gambar
%   - gaya Daftar Pustaka Amikom (biblatex)
%   - makro-makro bantu template
%============================================================

% ── Paragraf ─────────────────────────────────────────────────
\setlength{\parindent}{1.27cm}    % perataan awal paragraf (standar Word 0.5")
\setlength{\parskip}{0pt}         % tanpa jarak antar-paragraf
\clubpenalty=10000                % cegah baris yatim/piatu
\widowpenalty=10000

% Spasi utama naskah (Bab I–V, termasuk judul bab): spasi 2.
% Halaman-halaman awal mengatur spasinya sendiri di file masing-masing.
\doublespacing

% ── Penomoran ────────────────────────────────────────────────
% Bab di daftar isi & judul pakai angka Romawi: BAB I, BAB II, ...
% Tabel/gambar/rumus tetap arab: Tabel 2.1, Gambar 2.1, Persamaan (2.1).
\renewcommand{\thechapter}{\Roman{chapter}}
\renewcommand{\thetable}{\arabic{chapter}.\arabic{table}}
\renewcommand{\thefigure}{\arabic{chapter}.\arabic{figure}}
\renewcommand{\theequation}{\arabic{chapter}.\arabic{equation}}
% Sub-bab pakai huruf: A., B., C. (bukan 1.1, 1.2) — sesuai template asli.
\renewcommand{\thesection}{\Alph{section}}
\renewcommand{\thesubsection}{\arabic{subsection}}
\renewcommand{\thesubsubsection}{\alph{subsubsection}}

% ── Format judul bab ─────────────────────────────────────────
% Dua baris terpusat bold:
%       BAB I
%       PENDAHULUAN
% Spasi 2 tetap berlaku di judul bab (termasuk judul bab).
\titleformat{\chapter}[display]
  {\centering\normalsize\bfseries}
  {BAB \thechapter}
  {0.6em}
  {\MakeUppercase}
% Sebelum/sesudah judul bab dikontrol ketat agar tidak ada jarak berlebih.
\titlespacing*{\chapter}{0pt}{-1.2em}{0.6em}

% Sub-bab:  A. Latar Belakang Masalah   (bold, rata kiri)
\titleformat{\section}
  {\normalsize\bfseries}{\thesection.}{0.5em}{}
\titlespacing*{\section}{0pt}{0.4em}{0.2em}

% Sub-sub-bab:  1. Judul   (bold, rata kiri)
\titleformat{\subsection}
  {\normalsize\bfseries}{\thesubsection.}{0.5em}{}
\titlespacing*{\subsection}{0pt}{0.4em}{0.2em}

% Tingkat keempat:  a. Judul   (bold, rata kiri)
\titleformat{\subsubsection}
  {\normalsize\bfseries}{\thesubsubsection.}{0.5em}{}
\titlespacing*{\subsubsection}{0pt}{0.4em}{0.2em}

% ── Daftar isi / tabel / gambar (tocloft) ────────────────────
% Hanya BAB dan sub-bab yang masuk daftar isi (level 2), dan
% keduanya dicetak tebal. Teks daftar: Times New Roman 12,
% spasi 1,5, tidak di-bold (kecuali bab & sub-bab).
\setcounter{tocdepth}{1}
\setcounter{secnumdepth}{3}       % sub-sub-bab bernomor di naskah

\renewcommand{\cftchapfont}{\bfseries}
\renewcommand{\cftchappagefont}{\bfseries}
\renewcommand{\cftchappresnum}{BAB }
\renewcommand{\cftchapaftersnum}{\hspace{0.4em}}
\setlength{\cftchapnumwidth}{3.2em}        % cukup untuk "BAB IV"
\renewcommand{\cftchapleader}{\cftdotfill{\cftdotsep}}

\renewcommand{\cftsecfont}{\bfseries}
\renewcommand{\cftsecpagefont}{\bfseries}
\renewcommand{\cftsecfont}{\bfseries}
\renewcommand{\cftsecpagefont}{\bfseries}
\renewcommand{\cftsecaftersnum}{.}
\setlength{\cftsecnumwidth}{1.8em}         % lebar "A."
\setlength{\cftsecindent}{1.5em}
\setlength{\cftbeforechapskip}{0em}        % rapat agar daftar isi muat satu halaman
% Spasi judul halaman daftar (DAFTAR ISI/TABEL/GAMBAR/LAMPIRAN)
% dirapatkan agar daftar isi tidak meluber ke halaman berikutnya.
\setlength{\cftbeforetoctitleskip}{1em}
\setlength{\cftaftertoctitleskip}{1em}

% Entri DAFTAR TABEL / DAFTAR GAMBAR memakai prefix seperti
% caption-nya: "Tabel 2.1. Data mahasiswa" / "Gambar 2.1. Logo lama".
\renewcommand{\cfttabpresnum}{Tabel~}
\renewcommand{\cfttabaftersnum}{.}
\setlength{\cfttabnumwidth}{4.6em}
\renewcommand{\cftfigpresnum}{Gambar~}
\renewcommand{\cftfigaftersnum}{.}
\setlength{\cftfignumwidth}{4.6em}

% ── Caption tabel & gambar ───────────────────────────────────
% Tabel  : caption DI ATAS tabel, format "Tabel X.Y. Judul".
% Gambar : caption DI BAWAH gambar, format "Gambar X.Y. Judul".
\captionsetup{%
  labelsep=period,
  font={normalsize,stretch=1},
  justification=centering,
  singlelinecheck=true}
\captionsetup[table]{position=above,skip=6pt}
\captionsetup[figure]{position=below,skip=6pt}

% ── Nomor halaman ────────────────────────────────────────────
% Awal halaman bab & halaman bagian awal : bawah tengah (plain)
% Halaman lanjutan                        : kanan atas  (utama)
% DAFTAR PUSTAKA & LAMPIRAN               : tanpa nomor halaman
\setlength{\headheight}{14.5pt}       % cukup untuk fancyhdr
\fancypagestyle{plain}{%
  \fancyhf{}%
  \fancyfoot[C]{\thepage}%
  \renewcommand{\headrulewidth}{0pt}%
  \renewcommand{\footrulewidth}{0pt}%
}
\fancypagestyle{utama}{%
  \fancyhf{}%
  \fancyhead[R]{\thepage}%
  \renewcommand{\headrulewidth}{0pt}%
  \renewcommand{\footrulewidth}{0pt}%
}

% ── Nama daftar ──────────────────────────────────────────────
% babel-indonesian menimpa nama-nama ini saat \begin{document},
% jadi perubahan dibungkus \AtBeginDocument.
\AtBeginDocument{%
  \renewcommand{\contentsname}{DAFTAR ISI}%
  \renewcommand{\listtablename}{DAFTAR TABEL}%
  \renewcommand{\listfigurename}{DAFTAR GAMBAR}%
}

%============================================================
% MAKRO-MAKRO BANTU TEMPLATE
%============================================================

% \TeksTersembunyi{teks} — teks putih (tak terlihat di PDF).
% Template asli menaruh tulisan "HALAMAN SAMPUL", "HALAMAN JUDUL",
% dst. berwarna putih agar dapat terbaca oleh generator daftar isi.
\newcommand{\TeksTersembunyi}[1]{\textcolor{white}{#1}}

% \HalamanAwal{NAMA HALAMAN} — teks putih tersembunyi + masuk
% DAFTAR ISI (tanpa nomor bab, dengan nomor halaman romawi).
\makeatletter
\newcommand{\HalamanAwal}[1]{%
  \par\noindent\TeksTersembunyi{#1}\par
  \addcontentsline{toc}{halamanawal}{#1}%
}
% Entri daftar isi untuk halaman awal: normal (tidak tebal), ada titik-titik
% dan nomor halaman romawi.
\newcommand{\l@halamanawal}{\@dottedtocline{1}{0pt}{0pt}}

% \TocHalaman{NAMA} — entri DAFTAR ISI untuk halaman yang sudah
% punya judul terlihat (KATA PENGANTAR, DAFTAR ISI, DAFTAR TABEL,
% ABSTRAK, dan sebagainya). Tanpa teks putih tambahan.
\newcommand{\TocHalaman}[1]{\addcontentsline{toc}{halamanawal}{#1}}

% \BagianAkhir{NAMA BAGIAN} — entri DAFTAR PUSTAKA / LAMPIRAN:
% tebal, TANPA titik-titik dan TANPA nomor halaman (sesuai template asli).
\newcommand{\BagianAkhir}[1]{%
  \addtocontents{toc}{\protect\contentsline{bagianakhir}{\protect#1}{}{}}%
}
\newcommand{\l@bagianakhir}[2]{%
  \vskip 0.1em
  \noindent\textbf{#1}\par
}
\makeatother

% \Istilah{istilah}{penjelasan} — entri DAFTAR ISTILAH.
\newcommand{\Istilah}[2]{%
  \noindent\textbf{#1} : #2\par
}

% \Lampiran{Judul} — mulai lampiran baru: "Lampiran N. Judul"
% (bold), otomatis masuk DAFTAR LAMPIRAN.
\newcounter{lampiran}
\newcommand{\LampiranList}[1]{\noindent #1\par\vspace{0.4em}}
\makeatletter
\newcommand{\Lampiran}[1]{%
  \clearpage
  \stepcounter{lampiran}%
  \noindent\textbf{Lampiran \arabic{lampiran}. #1}\par\vspace{0.5em}
  \addtocontents{lampl}{\protect\LampiranList{Lampiran \arabic{lampiran}. #1}}%
}
% \listoflampiran — halaman DAFTAR LAMPIRAN (tanpa titik-titik
% dan tanpa nomor halaman pada entri-entrinya).
\newcommand{\listoflampiran}{%
  \chapter*{DAFTAR LAMPIRAN}%
  \TocHalaman{DAFTAR LAMPIRAN}%
  \@starttoc{lampl}%
}
\makeatother

% \KotakMaterai — kotak pengganti materai Rp10.000 pada Surat
% Pernyataan Keaslian Penelitian (diganti materai asli saat
% dicetak dan ditandatangani).
\newcommand{\KotakMaterai}{%
  \fbox{\parbox[c][2.8cm][c]{4.5cm}{\centering Bermaterai\\ Rp10.000}}%
}

% \KotakScan — kotak penanda tempat memasukkan hasil scan pada
% lampiran (kartu bimbingan, surat keterangan, dll.).
\newcommand{\KotakScan}[1]{%
  \noindent\fbox{\parbox[c][8cm][c]{\dimexpr\linewidth-2\fboxsep-2\fboxrule\relax}{%
    \centering #1}}%
}

% ── Lingkungan tabel (longtable) ─────────────────────────────
% \begin{tabel}{spesifikasi-kolom} ... \end{tabel}
% Isi tabel otomatis: Times New Roman 10, spasi 1, justify.
% Lihat contoh pemakaian lengkap (dengan tabel lanjutan) di
% chapters/bab2.tex.
\newenvironment{tabel}[1]{%
  \begingroup
  \FontTabel
  \setstretch{1.0}
  \setlength{\LTpre}{6pt}%
  \setlength{\LTpost}{6pt}%
  \begin{longtable}{#1}%
}{%
  \end{longtable}%
  \endgroup
}

% \TabelLanjutan{jumlah-kolom}{Judul Tabel} — kepala halaman
% lanjutan tabel: "Tabel X.Y. Judul Tabel (Lanjutan)" + kepala tabel.
% Pemakaian di dalam \endhead (lihat contoh chapters/bab2.tex).
\newcommand{\TabelLanjutan}[2]{%
  \multicolumn{#1}{c}{Tabel \thetable. #2 (Lanjutan)}\\
}

% ── Halaman landscape (kertas miring) ────────────────────────
% Nomor halaman pada halaman miring tetap terbaca horizontal di
% KANAN ATAS orientasi miring (sesuai contoh halaman landscape di
% template asli). Caranya: glyph dicetak di kiri-atas lembar fisik
% dengan rotasi 90°, sehingga setelah halaman dirotasi viewer
% (/Rotate) ia tampil horizontal di pojok kanan atas.
\fancypagestyle{landskape}{%
  \fancyhf{}%
  \fancyhead[L]{\rotatebox{90}{\thepage}}%
  \renewcommand{\headrulewidth}{0pt}%
  \renewcommand{\footrulewidth}{0pt}%
}
% Lingkungan landskape — pembungkus nyaman:
%   \begin{landskape} ... konten/tabel ... \end{landskape}
\newenvironment{landskape}{%
  \begin{landscape}%
  \pagestyle{landskape}%
}{%
  \end{landscape}%
  \pagestyle{utama}%
}

%============================================================
% DAFTAR PUSTAKA — GAYA AMIKOM (biblatex)
%============================================================
% Judul ditulis apa adanya (Title Case dari file .bib, tidak
% dirubah jadi sentence case oleh biblatex).
\DeclareFieldFormat{titlecase}{#1}

% Pengarang terakhir dihubungkan "dan" (bukan "&"):
%   Wibowo, A. R., dan Saputra, I. A.
\DeclareDelimFormat[bib,biblist]{finalnamedelim}{\addcomma\space dan}

% Kata-kata lokal ("dan", "dkk.", "Diambil dari", dsb.) disediakan
% oleh indonesian.lbx — file lokalisasi biblatex milik repo,
% karena biblatex belum menyediakan indonesian.lbx.
% Judul artikel/tesis tanpa tanda kutip (gaya Amikom: judul polos,
% yang italic hanya nama jurnal/judul buku).
\DeclareFieldFormat[article,inbook,incollection,inproceedings,patent,thesis,unpublished]{title}{#1\isdot}
% Volume(nomor) halaman tanpa awalan "pp." : "Jurnal Teknik Elektro, 11(1), 26-30"
\DeclareFieldFormat{pages}{#1}
\renewbibmacro*{journal+issuetitle}{%
  \usebibmacro{journal}%
  \setunit*{\addspace}%
  \printfield[journaltitle]{series}%
  \setunit{\addcomma\space}%
  \printfield[nq]{issue}%
  \setunit{\addcomma\space}%
  \usebibmacro{volume+number+eid}%
  \setunit{\addcomma\space}%
  \usebibmacro{issue+date}%
  \setunit{\addcolon\space}%
  \usebibmacro{issue}%
  \newunit}
\renewbibmacro*{volume+number+eid}{%
  \printfield{volume}%
  \iffieldundef{number}{}{\mkbibparens{\printfield{number}}}%
  \setunit{\addcomma\space}%
  \printfield{eid}%
}
% Jenis skripsi: "Skripsi." (bukan "Master's thesis.")
\DeclareFieldFormat[thesis]{type}{#1\adddot}

% Hilangkan awalan "in:" pada prosiding (gaya Amikom tidak memakainya)
\renewbibmacro{in:}{}

% Jarak penulisan: dalam referensi spasi 1, antar referensi 2 spasi.
\renewcommand{\bibfont}{\setstretch{1.0}}
\setlength{\bibitemsep}{2\baselineskip}
\setlength{\bibhang}{1.27cm}
